% -*- coding: utf-8 -*-
% !TEX program = xelatex
\documentclass[plain,noanswer,medmath]{../../template/jnuexam}
\usepackage{../../template/kysx}

\begin{document}


\examtitle{2002}{1}


\makepart{填空题}{本题共5小题，每小题3分，满分15分}


\begin{problem}
$\int_{\e}^{+\infty} \frac{\dx}{x\ln^2 x} =$ \fillin{}．
\end{problem}


\begin{problem}
已知函数$y = y(x)$由方程$\e^y + 6xy + x^2 - 1 = 0$确定，则$y''(0) =$ \fillin{}．
\end{problem}


\begin{problem}
微分方程$yy'' + (y')^2 = 0$满足初始条件$y|_{x = 0} = 1$, $y' |_{x = 0} = \frac{1}{2}$的特解是 \fillin{}．
\end{problem}


\begin{problem}
已知实二次型$f(x_1,x_2,x_3) = a(x_1^2 + x_2^2 + x_3^2) + 4x_1x_2 + 4x_1x_3 + 4x_2x_3$
经正交变换$x = Py$可化成标准型$f = 6y_1^2$，则$a =$ \fillin{}．
\end{problem}


\begin{problem}
设随机变量$X$服从正态分布$N(\mu,\sigma^2)$（$\sigma > 0$），
且二次方程$y^2 + 4y + X = 0$无实根的概率为$\frac{1}{2}$，则$\mu =$ \fillin{}．
\end{problem}


\makepart{选择题}{本题共5小题，每小题3分，共15分}


\begin{problem}
考虑二元函数$f(x,y)$的下面4条性质：\par
①$f(x,y)$在点$(x_0,y_0)$处连续，	②$f(x,y)$在点$(x_0,y_0)$处的两个偏导数连续，\par
③$f(x,y)$在点$(x_0,y_0)$处可微，	④$f(x,y)$在点$(x_0,y_0)$处的两个偏导数存在．\par
若用``$P \Rightarrow Q$''表示可由性质$P$推出$Q$，则有\pickout{}
\begin{abcd}
\item ②$\Rightarrow$③$\Rightarrow$①．
\item ③$\Rightarrow$②$\Rightarrow$①．
\item ③$\Rightarrow$④$\Rightarrow$①．
\item ③$\Rightarrow$①$\Rightarrow$④．
\end{abcd}
\end{problem}


\begin{problem}
设$u_n \ne 0$（$n = 1,2,3, \cdots$），且$\lim_{n \to \infty} \frac{n}{u_n} = 1$，
则级数$\sum_{n = 1}^{\infty}(- 1)^{n + 1} \left(\frac{1}{u_n} + \frac{1}{u_{n + 1}} \right)$\pickout{}
\begin{abcd}
\item 发散．
\item 绝对收敛．
\item 条件收敛．
\item 收敛性根据所给条件不能判定．
\end{abcd}
\end{problem}


\begin{problem}
设函数$y = f(x)$在$(0, +\infty)$内有界且可导，则\pickout{}
\begin{abcd}
\item 当$\lim_{x \to +\infty} f(x) = 0$时，必有$\lim_{x \to +\infty} f'(x) = 0$．
\item 当$\lim_{x \to +\infty} f'(x)$存在时，必有$\lim_{x \to +\infty} f'(x) = 0$．
\item 当$\lim_{x \to 0^+} f(x) = 0$时，必有$\lim_{x \to 0^+} f'(x) = 0$．
\item 当$\lim_{x \to 0^+} f'(x)$存在时，必有$\lim_{x \to 0^+} f'(x) = 0$．
\end{abcd}
\end{problem}


\begin{problem}
设有三张不同平面的方程$a_{i1} x + a_{i2} y + a_{i3} z = b_i,i = 1,2,3$，
它们所组成的线性方程组的系数矩阵与增广矩阵的秩都为$2$，则这三张平面可能的位置关系为\pickout{}
\begin{abcd}
\item \includegraphics[width=0.18\linewidth]{2002A1}
\item 
\includegraphics[width=0.18\linewidth]{2002A2}
\item \includegraphics[width=0.18\linewidth]{2002A3}
\item \includegraphics[width=0.18\linewidth]{2002A4}
\end{abcd}
\end{problem}


\begin{problem}
设$X_1$和$X_2$是任意两个相互独立的连续型随机变量，它们的概率密度分别为$f_1(x)$和$f_2(x)$，分布函数分别为$F_1(x)$和$F_2(x)$，则\pickout{}
\begin{abcd}
\item $f_1(x) + f_2(x)$必为某一随机变量的概率密度．
\item $f_1(x) f_2(x)$必为某一随机变量的概率密度．
\item $F_1(x) + F_2(x)$必为某一随机变量的分布函数．
\item $F_1(x) F_2(x)$必为某一随机变量的分布函数．
\end{abcd}
\end{problem}


\makepart{}{本题满分6分}

\begin{problem}
设函数$f(x)$在$x = 0$的某邻域内具有一阶连续导数，且$f(0) \ne 0,f'(0) \ne 0$，
若$af(h) + bf(2h) - f(0)$在$h \to 0$时是比$h$高阶的无穷小，试确定$a,b$的值．
\end{problem}


\makepart{}{本题满分7分}

\begin{problem}
已知两曲线$y = f(x)$与$y = \int_0^{\arctan x} \!\e^{- t^2} \dt$在点$(0,0)$处的切线相同，
写出此切线方程，并求极限$\lim_{n \to \infty} nf\left(\frac{2}{n} \right)$．
\end{problem}


\makepart{}{本题满分7分}

\begin{problem}
计算二重积分$\iint_D \e^{\max \{x^2,y^2\}} \dx\dy$，其中$D = \{(x,y)|0 \le x \le 1,0 \le y \le 1\}$．
\end{problem}


\makepart{}{本题满分8分}

\begin{problem}
设函数$f(x)$在$(-\infty , +\infty)$内具有一阶连续导数，$L$是上半平面$(y > 0)$内的有向分段光滑曲线，
其起点为$(a,b)$，终点为$(c,d)$．记
$$I = \int_L \frac{1}{y} [1 + y^2 f(xy)]\dx + \frac{x}{y^2}[y^2 f(xy) - 1]\dy,$$
\begin{enumerate*}
\item 证明曲线积分$I$与路径$L$无关；\item 当$ab = cd$时，求$I$的值．
\end{enumerate*}

\end{problem}


\makepart{}{本题满分7分}

\begin{problem}
\begin{enumerate}
  \item 验证函数$y(x) = 1 + \frac{x^3}{3!} + \frac{x^6}{6!} + \frac{x^9}{9!}
        + \cdots + \frac{x^{3n}}{(3n)!}+ \cdots$（$-\infty < x < +\infty$）满足微分方程$y'' + y' + y = \e^x$；
  \item 利用(1)的结果求幂级数$\sum_{n = 0}^{\infty}\frac{x^{3n}}{(3n)!}$的和函数．
\end{enumerate}
\end{problem}


\makepart{}{本题满分7分}

\begin{problem}
设有一小山，取它的底面所在的平面为$xoy$坐标面，其底部所占的区域为
$D = \left\{(x,y)\left| x^2 + y^2 - xy \le 75 \right. \right\}$，小山的高度函数为$h(x,y) = 75 - x^2 - y^2 + xy$．
\begin{enumerate}
\item 设$M(x_0,y_0)$为区域$D$上的一点，问$h(x,y)$在该点沿平面上什么方向的方向导数最大？
若记此反向导数的最大值为$g(x_0,y_0)$，试写出$g(x_0,y_0)$表达式．
\item 现欲利用此小山开展攀岩活动，为此需要在山脚寻找一上山坡度最大的点作为攀登的起点．也就是说，
要在$D$的边界线$x^2 + y^2 - xy = 75$上找出使(\!Ⅰ\!)中的$g(x,y)$达到最大值的点．试确定攀登起点的位置．
\end{enumerate}
\end{problem}


\makepart{}{本题满分6分}

\begin{problem}
已知$4$阶方阵$A = (\alpha_1,\alpha_2,\alpha_3,\alpha_4)$，$\alpha_1,\alpha_2,\alpha_3,\alpha_4$均为4维列向量，
其中$\alpha_2,\alpha_3,\alpha_4$线性无关，$\alpha_1 = 2\alpha_2 - \alpha_3$．
如果$\beta = \alpha_1 + \alpha_2 + \alpha_3 + \alpha_4$，求线性方程组$Ax = \beta$的通解．
\end{problem}


\makepart{}{本题满分8分}

\begin{problem}
设$A,B$为同阶方阵，
\begin{enumerate}
\item 如果$A,B$相似，试证$A,B$的特征多项式相等．
\item 举一个二阶方阵的例子说明(\!Ⅰ\!)的逆命题不成立．
\item 当$A,B$均为实对称矩阵时，试证(\!Ⅰ\!)的逆命题成立．
\end{enumerate}
\end{problem}


\makepart{}{本题满分8分}

\begin{problem}
设随机变量$X$的概率密度为$f(x) = \begin{cases}
  \frac{1}{2}\cos \frac{x}{2} & 0 \le x \le \pi ; \\
                           0, & \text{其他}.
\end{cases}$对$X$独立地重复观察$4$次，用$Y$表示观察值大于$\frac{\pi}{3}$的次数，求$Y^2$的数学期望．
\end{problem}


\makepart{}{本题满分8分}

\begin{problem}
设总体$X$的概率分布为
$$\begin{array}{|c|cccc|}
\hline
  X	& 0	       & 1                 & 2        & 3 \\
\hline
  P	& \theta^2 & 2\theta(1-\theta) & \theta^2 &	1-2\theta \\
\hline
\end{array}$$
其中$\theta$（$\text{0 <}\theta \text{<}\frac{1}{2}$）是未知参数，
利用总体$X$的如下样本值$3,1,3,0,3,1,2,3$，求$\theta$的矩阵估计值和最大似然估计值．
\end{problem}


\end{document}
